% GENERATED FILE. Edit canonical lessons, then run npm run generate:books.
% canonical-source-hash: fnv1a64:894de2d3d726eb58
% canonical-lessons: ES-C556-entierro, ES-C556-enterrar, ES-C556-velatorio, ES-C556-tumba, ES-C556-cementerio

\chapter{El entierro}
\label{ch:es-el-entierro}
\hlchaptermodality{556}

\begin{chapteropening}
\textbf{By the end of this chapter:} \emph{I can talk about a funeral and a burial in a cemetery.}

Complete the last lesson of chapter 556: I can talk about a funeral and a burial in a cemetery.
\end{chapteropening}

\section[el entierro]{el entierro — the burial, the funeral}

\label{lesson:ES-C556-entierro}



\begin{quote}
Before the new one: say the Spanish for dead, then the Spanish for the deceased.
\end{quote}

\subsection*{You'll want to know: el entierro}
\textbf{el entierro} — \textquotedblleft{}the burial, the funeral\textquotedblright{}.

\begin{quote}
Fuimos al entierro de su abuelo.
\end{quote}

\begin{quote}
\emph{We went to his grandfather's funeral.}
\end{quote}

\begin{cousinweb}
\emph{entierro} is built on \emph{enterrar} 'to bury', which comes from spoken Latin \emph{*interrāre} 'to put in the earth' (EtymDB), from \emph{terra} 'earth'.
\end{cousinweb}

\begin{grammarlens}[title={What you've built}]
One more word for telling, explaining and describing at B1.
\end{grammarlens}

\subsection*{Your turn}
\begin{itemize}
\raggedright
  \item \emph{Say it:} \emph{el entierro}
  \item \emph{Say it:} \emph{el entierro}, once more
  \item \emph{Say it:} a sentence of your own with \emph{el entierro}
  \item \emph{Recall:} say the Spanish for dead, then the Spanish for the deceased, then say \emph{el entierro} again
\end{itemize}

\begin{tcolorbox}[breakable,colback=teal!4,colframe=teal!35!black,title={Before you move on}]
What does el entierro mean? (\textquotedblleft{}The burial, the funeral\textquotedblright{}.) Say it once more.
\end{tcolorbox}
\section[enterrar]{enterrar — to bury}

\label{lesson:ES-C556-enterrar}



\begin{quote}
Before the new one: say the Spanish for the deceased, then the Spanish for the burial.
\end{quote}

\subsection*{You'll want to know: enterrar}
\textbf{enterrar} — \textquotedblleft{}to bury\textquotedblright{}.

\begin{quote}
Lo enterraron en su pueblo.
\end{quote}

\begin{quote}
\emph{They buried him in his home village.}
\end{quote}

Stem change: \emph{entierro}, \emph{entierras}.

\begin{cousinweb}
\emph{enterrar} is inherited from spoken Latin \emph{*interrāre} 'to put in the earth' (EtymDB), built on \emph{in} and \emph{terra} 'earth'.
\end{cousinweb}

\begin{grammarlens}[title={What you've built}]
One more word for telling, explaining and describing at B1.
\end{grammarlens}

\subsection*{Your turn}
\begin{itemize}
\raggedright
  \item \emph{Say it:} \emph{enterrar}
  \item \emph{Say it:} \emph{enterrar}, once more
  \item \emph{Say it:} a sentence of your own with \emph{enterrar}
  \item \emph{Recall:} say the Spanish for the deceased, then the Spanish for the burial, then say \emph{enterrar} again
\end{itemize}

\begin{tcolorbox}[breakable,colback=teal!4,colframe=teal!35!black,title={Before you move on}]
What does enterrar mean? (\textquotedblleft{}To bury\textquotedblright{}.) Say it once more.
\end{tcolorbox}
\section[el velatorio]{el velatorio — the wake; the funeral parlour}

\label{lesson:ES-C556-velatorio}



\begin{quote}
Before the new one: say the Spanish for the burial, then the Spanish for to bury.
\end{quote}

\subsection*{You'll want to know: el velatorio}
\textbf{el velatorio} — \textquotedblleft{}the wake; the funeral parlour\textquotedblright{}.

\begin{quote}
Fuimos al velatorio por la noche.
\end{quote}

\begin{quote}
\emph{We went to the wake in the evening.}
\end{quote}

In Latin America \emph{el velorio} is common.

\begin{cousinweb}
\emph{velatorio} is built on \emph{velar} 'to stay awake, to keep watch over', which is inherited from Latin \emph{vigilāre} 'to stay awake' (EtymDB). Family and friends stay up with the dead.
\end{cousinweb}

\begin{grammarlens}[title={What you've built}]
One more word for telling, explaining and describing at B1.
\end{grammarlens}

\subsection*{Your turn}
\begin{itemize}
\raggedright
  \item \emph{Say it:} \emph{el velatorio}
  \item \emph{Say it:} \emph{el velatorio}, once more
  \item \emph{Say it:} a sentence of your own with \emph{el velatorio}
  \item \emph{Recall:} say the Spanish for the burial, then the Spanish for to bury, then say \emph{el velatorio} again
\end{itemize}

\begin{tcolorbox}[breakable,colback=teal!4,colframe=teal!35!black,title={Before you move on}]
What does el velatorio mean? (\textquotedblleft{}The wake; the funeral parlour\textquotedblright{}.) Say it once more.
\end{tcolorbox}
\section[la tumba]{la tumba — the grave, the tomb}

\label{lesson:ES-C556-tumba}



\begin{quote}
Before the new one: say the Spanish for to bury, then the Spanish for the wake.
\end{quote}

\subsection*{You'll want to know: la tumba}
\textbf{la tumba} — \textquotedblleft{}the grave, the tomb\textquotedblright{}.

\begin{quote}
Pusimos flores en su tumba.
\end{quote}

\begin{quote}
\emph{We put flowers on his grave.}
\end{quote}

\begin{cousinweb}
\emph{tumba} comes from Late Latin \emph{tumba}, taken from Greek \emph{týmbos} 'a burial mound' (EtymDB).
\end{cousinweb}

\begin{grammarlens}[title={What you've built}]
One more word for telling, explaining and describing at B1.
\end{grammarlens}

\subsection*{Your turn}
\begin{itemize}
\raggedright
  \item \emph{Say it:} \emph{la tumba}
  \item \emph{Say it:} \emph{la tumba}, once more
  \item \emph{Say it:} a sentence of your own with \emph{la tumba}
  \item \emph{Recall:} say the Spanish for to bury, then the Spanish for the wake, then say \emph{la tumba} again
\end{itemize}

\begin{tcolorbox}[breakable,colback=teal!4,colframe=teal!35!black,title={Before you move on}]
What does la tumba mean? (\textquotedblleft{}The grave, the tomb\textquotedblright{}.) Say it once more.
\end{tcolorbox}
\section[el cementerio]{el cementerio — the cemetery}

\label{lesson:ES-C556-cementerio}



\begin{quote}
Before the new one: say the Spanish for the wake, then the Spanish for the grave.
\end{quote}

\subsection*{You'll want to know: el cementerio}
\textbf{el cementerio} — \textquotedblleft{}the cemetery\textquotedblright{}.

\begin{quote}
El cementerio está detrás de la iglesia.
\end{quote}

\begin{quote}
\emph{The cemetery is behind the church.}
\end{quote}

\begin{cousinweb}
\emph{cementerio} is a learned borrowing from Church Latin \emph{coemētērium} (EtymDB), taken from Greek \emph{koimēt\'{\={e}}rion} 'a sleeping place'.
\end{cousinweb}

\begin{grammarlens}[title={What you've built}]
One more word for telling, explaining and describing at B1.
\end{grammarlens}

\subsection*{Your turn}
\begin{itemize}
\raggedright
  \item \emph{Say it:} \emph{el cementerio}
  \item \emph{Say it:} \emph{el cementerio}, once more
  \item \emph{Say it:} a sentence of your own with \emph{el cementerio}
  \item \emph{Recall:} say the Spanish for the wake, then the Spanish for the grave, then say \emph{el cementerio} again
\end{itemize}

\begin{tcolorbox}[breakable,colback=teal!4,colframe=teal!35!black,title={Before you move on}]
What does el cementerio mean? (\textquotedblleft{}The cemetery\textquotedblright{}.) Say it once more.
\end{tcolorbox}
